% ==================== 第 3 章 研究目标与主要内容 ====================
\clearpage
\section{研究目标、主要内容与关键问题}

\subsection{总体目标}

本课题拟构建一套零硬冲突、可解释、可复核、可扩展的智能排课决策系统。
系统接收结构化教师、教室、班级、课程和时间槽数据，输出满足全部硬约束的
课表；在存在多个可行方案时，根据不同业务策略优化软目标，并提供评分、
原因说明和变更明细。最终形成可独立运行的原型，同时完成与飞书协同流程的
接口设计和验证方案。

\subsection{主要研究内容}

\begin{enumerate}
  \item \textbf{业务对象与约束建模。}明确五类核心对象及其关系，把资质、
  容量、设备、不可用时段和资源互斥转化为可计算约束。
  \item \textbf{多目标排课模型。}在硬约束之外，引入学生偏好、教师偏好、
  容量适配、时间紧凑和负载均衡等软目标，并设计可配置策略。
  \item \textbf{最小影响调课。}以原课表为基线，把变更课次数量纳入目标，
  生成新课表及逐项差异。
  \item \textbf{校验与解释。}建立独立校验器，复核课次完整性和资源冲突；
  设计分项指标，使教务人员能够理解方案取舍。
  \item \textbf{应用与平台集成。}实现本地 API 和交互界面，设计飞书多维
  表格数据映射、消息卡片和审批回调方案。
  \item \textbf{实验与评价。}围绕正确性、性能、稳定性、偏好敏感性和
  可用性开展实验，形成可复现的评价流程。
\end{enumerate}

\subsection{拟解决的关键问题}

\subsubsection{硬约束的完备表达}

需要保证所有实际不可违反的规则都进入模型，并避免同一规则在前端、后端和
数据表中以不同口径重复实现。课题将以统一数据模型为入口，以求解模型和
独立校验器分别实现约束检查，降低漏约束风险。

\subsubsection{多目标之间的可解释权衡}

学生偏好、教师偏好和资源效率可能相互冲突。课题不直接比较不同权重体系下
的 CP-SAT 原始目标值，而采用统一口径的分项指标和归一化评分比较方案；
同时通过权重敏感性实验分析方案变化。

\subsubsection{动态变化下的方案稳定性}

调课不应简单地重新生成完全不同的课表。系统需要识别受影响课次，以原方案
为基线，在满足新增硬约束的前提下最小化变更数量，并输出变更原因和影响对象。

\subsection{技术指标与验收方式}

表 \ref{tab:targets} 中的数值是本课题拟达到的验收目标，不是外部行业统计。

\begin{table}[htbp]
  \centering
  \caption{研究目标与验收方式}
  \label{tab:targets}
  \small
  \begin{tabularx}{\textwidth}{L{2.6cm} L{3.2cm} Y Y}
    \toprule
    \textbf{维度} & \textbf{目标} & \textbf{验收方法} & \textbf{证据形式} \\
    \midrule
    正确性 & 可行方案硬冲突数为 0 & 求解后运行独立校验器，并覆盖不可行场景 & 自动化测试与校验报告 \\
    完整性 & 所有课程请求按规定课次排入 & 对课程请求与输出课次逐项核对 & 课次清单与缺失检查 \\
    性能 & 典型规模在设定时限内返回可行结果 & 分层构造不同规模数据并重复测量 & 运行环境、实例规模与耗时记录 \\
    稳定性 & 调课时优先减少相对基线的变更 & 设置教师不可用事件并统计变更课次 & 调课前后差异清单 \\
    可解释性 & 展示至少三类分项指标及方案差异 & 检查评分卡、解释文本和权重配置 & 界面截图与指标定义 \\
    可集成性 & 提供结构化接口并完成飞书映射设计 & 验证 JSON 输入输出及回调流程 & API 文档与集成原型 \\
    \bottomrule
  \end{tabularx}
\end{table}
